\chapter*{Nomenclatura}
\addcontentsline{toc}{chapter}{Nomenclatura}

\section*{Abreviatures i Acrònims}

\begin{longtable}{p{3.2cm}p{10cm}}
    \toprule
    \textbf{Abreviatura} & \textbf{Definició} \\
    \midrule\endhead
    API      & \textit{Application Programming Interface} --- Interfície de programació d'aplicacions que permet la comunicació entre sistemes de programari. \\[4pt]
    CSS      & \textit{Cascading Style Sheets} --- Llenguatge d'estils per a la presentació visual de documents HTML. \\[4pt]
    DNS      & \textit{Domain Name System} --- Sistema de noms de domini que tradueix noms llegibles per humans en adreces IP. \\[4pt]
    HTML     & \textit{HyperText Markup Language} --- Llenguatge de marques estàndard per a la creació de pàgines web. \\[4pt]
    HTTPS    & \textit{HyperText Transfer Protocol Secure} --- Versió xifrada del protocol HTTP que garanteix la seguretat de les comunicacions web. \\[4pt]
    IEC      & Institut d'Estudis Catalans --- Organisme normatiu de la llengua catalana. \\[4pt]
    JSON     & \textit{JavaScript Object Notation} --- Format lleuger d'intercanvi de dades estructurades, llegible per humans i màquines. \\[4pt]
    REST     & \textit{Representational State Transfer} --- Estil d'arquitectura de programari per a serveis web. \\[4pt]
    TLS      & \textit{Transport Layer Security} --- Protocol criptogràfic que garanteix la seguretat de les comunicacions en xarxa. \\[4pt]
    UI       & \textit{User Interface} --- Interfície d'usuari; conjunt d'elements visuals i interactius amb els quals l'usuari interactua. \\[4pt]
    UIB      & Universitat de les Illes Balears. \\[4pt]
    UX       & \textit{User Experience} --- Experiència d'usuari; percepció global que té l'usuari en interaccionar amb un sistema. \\[4pt]
    WCAG     & \textit{Web Content Accessibility Guidelines} --- Directrius d'accessibilitat per al contingut web, publicades pel W3C. \\[4pt]
    W3C      & \textit{World Wide Web Consortium} --- Organisme internacional que desenvolupa estàndards per a la World Wide Web. \\[4pt]
    Schema.org & Iniciativa conjunta de Google, Bing i Yahoo! que defineix un vocabulari de metadades estructurades per a la web semàntica, facilitant la indexació semàntica i la generació de \textit{rich results} als motors de cerca. \\[4pt]
    SEO      & \textit{Search Engine Optimisation} --- Conjunt de tècniques per millorar la visibilitat i el posicionament d'un lloc web als resultats orgànics dels motors de cerca. \\[4pt]
    \bottomrule
\end{longtable}
